% RESEARCH_DRAFT_SHA256: 2d7a8f870a088feb68907e7ce2333fc2ca44a4cb86596dea25e081d1f5bedab9
% Layout and source-reference wrappers expanded for the web renderer.
\documentclass[11pt]{article}
\usepackage[margin=25mm]{geometry}
\usepackage{fontspec,unicode-math}
\setmainfont{Latin Modern Roman}
\setmathfont{Latin Modern Math}
\usepackage{amsmath,amssymb,mathtools,xurl,hyperref}
\hypersetup{colorlinks=true,urlcolor=blue,linkcolor=blue}
\setlength{\parindent}{0pt}
\setlength{\parskip}{5pt}
\setlength{\emergencystretch}{3em}
\newcommand{\sref}[2]{\hyperlink{source:#1:#2}{[S#2]}}
\newcommand{\cataloguescope}[1]{\paragraph{Source question.}#1}
\newcommand{\floor}[1]{\left\lfloor#1\right\rfloor}
\begin{document}

\section*{Multiplicity obstructions: an explicit small calculation and the orbit-closure gap}
\noindent UnsolvedMath ID 30004422; OWR-17472-018. Research attempt dated 11 October 2026.
\paragraph{Result in scope.}This manuscript contains a source-specific attempt and complete proofs of the stated auxiliary calculations. It makes no new solution or novelty claim for the entire target.
\section{Original problem and preserved definition}
The following definition and source list are reproduced from the repository. Any earlier inspection dates and reported literature claims in this preserved text describe that source record, not additional verification by this attempt.
% BEGIN PRESERVED DEFINITION
\subsection{Definitions and mathematical statement}

\paragraph{Shared notation.}Standard number systems are $\mathbb N,\mathbb Z,\mathbb Q,\mathbb R,\mathbb C$. Specialized notation and index ranges follow the cited source and the definitions below.


In the usual GCT setting, the general linear group $G=\operatorname{GL}(V)$ acts on a space of polynomials, with orbit $Gp$ and Zariski closure $\overline{Gp}$. Its coordinate ring is graded by polynomial degree. For an irreducible polynomial representation indexed by a partition $\lambda$, $\operatorname{mult}_\lambda\mathbb C[\overline{Gp}]_d$ is its multiplicity in the degree-$d$ part. Inclusion $\overline{Gp}\subseteq\overline{Gq}$ induces a surjection of coordinate rings in the reverse direction, so a larger multiplicity on the left obstructs such inclusion.
\paragraph{Record scope and qualifications.}
The catalogue title is retained as metadata. Named certificates, contracts, bounds or frameworks in that title are not asserted as results of the original source extract.
This is a family of multiplicity-obstruction questions. The supplied record does not specify $G$, $p$, $q$, or the size parameters; an answer must identify the particular geometric-complexity-theory family to which it applies.

Source-record limitation: The acting group and source polynomials p,q whose orbit closures are to be separated are not supplied; the core obstruction instance is undefined.
\paragraph{Statement recorded in the supplied source.}
For the polynomial orbit closures relevant to geometric complexity theory, can one find a partition $\lambda$ and degree $d$ such that
$$
\operatorname{mult}_\lambda\mathbb C[\overline{Gp}]_d>\operatorname{mult}_\lambda\mathbb C[\overline{Gq}]_d,
$$
thereby obtaining a multiplicity obstruction separating the orbit closures? \href{https://doi.org/10.4171/owr/2020/9}{[S2]}
\paragraph{Source status.}
The supplied catalogue labels this source record open; that metadata was checked against the supplied snapshot on 2026-10-11. The label is preserved for provenance and is not an independent verification of present open status or mathematical correctness. The catalogue literature assessment is dated 2026-08-22.
\subsection{Short English statement}
Separate polynomial orbit closures by finding an irreducible representation occurring more often in one coordinate ring than in the other.
\subsection{Sources}
\begin{itemize}
\item[S1] Ulam AI and the UnsolvedMath Contributors, UnsolvedMath: A Curated Collection of Open Mathematics Problems, record 30004422 (OWR-17472-018). Catalogue curation and metadata: CC BY 4.0. The source extract is retained; definitions, English summaries and qualifications are editorial additions.
\url{https://huggingface.co/datasets/ulamai/UnsolvedMath}.
\item[S2] Mini-Workshop: Kronecker, Plethysm, and Sylow Branching Coefficients and their Applications to Complexity Theory (2021), Oberwolfach Reports, DOI 10.4171/owr/2020/9, pp. 517--568
\url{https://doi.org/10.4171/owr/2020/9}.
\item[S3] Mulmuley and Sohoni, Geometric Complexity Theory I
\url{https://doi.org/10.1137/S009753970038715X}.
\item[S4] Ikenmeyer and Panova, Geometric complexity theory and matrix powering
\url{https://doi.org/10.1016/j.difgeo.2017.07.001}.
\item Source-specific Oberwolfach Reports licences, including CC BY-SA 4.0 where applicable, are retained; see THIRD\_PARTY\_NOTICES.md.
\end{itemize}
% END PRESERVED DEFINITION
\subsection{Verbatim supplied catalogue question}
For the polynomial orbit closures relevant to geometric complexity theory, can one find a partition $\lambda$ and degree $d$ such that
$$
\operatorname{mult}_\lambda\mathbb C[\overline{Gp}]_d>\operatorname{mult}_\lambda\mathbb C[\overline{Gq}]_d,
$$
thereby obtaining a multiplicity obstruction separating the orbit closures?

\section{What the source actually specifies}
The source report, printed pages 557--560, discusses padded permanent
versus determinant, then reports known multiplicity obstructions for
power sums versus products. Its concrete latter family is
$p=x_1^D+\cdots+x_m^D$, $q=x_1\cdots x_D$, with $G=GL_m$.
The supplied catalogue question omits which family and parameters
must be separated. We do not replace the hard permanent/determinant
problem by an easier unspecified pair and call that a solution.
The primary Ikenmeyer--Kandasamy abstract confirms that their reported
power-sum/product result already uses multiplicities beyond occurrence
obstructions. We have not reproduced its full representation-theoretic proof.

We attempt two elementary routes: first compute a complete small
occurrence obstruction to expose the required conventions, then test
whether elementary orbit separation or orbit multiplicities transfer
to the intended orbit closures.

\section{A full binary-quadratic auxiliary calculation}
This subsection uses a deliberately different, explicitly labeled
pair. Let $V$ be the space of complex symmetric matrices
$Q=\left(\begin{smallmatrix}a&b\\b&c\end{smallmatrix}\right)$,
equivalently binary quadratics $ax^2+2bxy+cy^2$. Let $GL_2$ act
by $Q\mapsto gQg^T$. Take $p=I$ and
$q=\operatorname{diag}(1,0)$.

The orbit closure of $p$ is all of $V$. To see this, every invertible
complex symmetric matrix is congruent to $I$. Indeed its symmetric
bilinear form has a vector with nonzero self-pairing: otherwise
$B(u+v,u+v)-B(u,u)-B(v,v)=2B(u,v)$ makes the whole form zero.
Scale that vector to self-pairing one using a complex square root,
take its orthogonal complement, and repeat in dimension one. This
constructs the congruence for an invertible form. Invertible symmetric
matrices are dense, because $Q+tI$ is invertible except at the
finitely many roots of its monic determinant polynomial. A polynomial
vanishing on the orbit consequently vanishes on all $V$ by continuity.

The orbit closure of $q$ is $Y=\{ac-b^2=0\}$. Every nonzero
rank-one symmetric form is congruent to $q$, by the same orthogonal
decomposition on its nonzero part, and zero is a limit under scalar
matrices. The displayed determinant equation defines a closed set
containing that orbit and no rank-two form, proving the claim.

Coordinate functions carry the contragredient action
$(g\cdot F)(Q)=F(g^{-1}Qg^{-T})$. The polynomial
$\Delta=ac-b^2$ of coordinate degree two spans the character
$g\mapsto\det(g)^{-2}$. In the conventional dual Schur notation
this is the dual of the partition $\lambda=(2,2)$; we spell out
the character to avoid a hidden sign convention for coordinate rings.

Its multiplicity in $\mathbb C[V]_2$ is exactly one. A diagonal
$g=\operatorname{diag}(t,s)$ acts on $a,b,c$ by factors
$t^{-2},t^{-1}s^{-1},s^{-2}$. Therefore a degree-two polynomial
with weight $t^{-2}s^{-2}$ is a combination $Aac+Bb^2$.
For $g=\left(\begin{smallmatrix}1&z\\0&1\end{smallmatrix}\right)$,
substitution sends $(a,b,c)$ to $(a-2zb+z^2c,b-zc,c)$.
Invariance under this determinant-one matrix forces $A+B=0$ by
the coefficients of $zbc$ and $z^2c^2$. Conversely $\Delta$
has the required character under every $g$, by multiplicativity
of determinant. The character space is precisely its one-dimensional span.

In $\mathbb C[Y]_2$ the same multiplicity is zero. The degree-two
relations vanishing on $Y$ are exactly multiples of $ac-b^2$:
substitute $(a,b,c)=(u^2,uv,v^2)$ into the six quadratic monomials.
Their images are $u^4,u^3v,u^2v^2,u^2v^2,uv^3,v^4$, with just
the indicated one-dimensional relation. The weight $t^{-2}s^{-2}$
subspace of the quotient is therefore spanned by the common class
of $ac$ and $b^2$. But the unipotent substitution sends the class
of $ac$ to that of $ac-2zbc+z^2c^2$, which is not fixed for all
$z$; the classes of $bc$ and $c^2$ are nonzero and independent
by the same substitution test. Thus no nonzero vector in that
weight space has the character, proving the multiplicity zero.
This gives the fully checked auxiliary pair $(\lambda,d)=((2,2),2)$
with multiplicities one and zero. It is a simple known occurrence
obstruction, not the source's nonoccurrence multiplicity construction.

\section{Testing the source power-sum/product route directly}
For the concrete source case $m=D=3$, the Fermat cubic
$p=x^3+y^3+z^3$ cannot be a product of three homogeneous linear
forms. Its gradient $(3x^2,3y^2,3z^2)$ vanishes only at zero.
If a cubic is a product $l_1l_2l_3$, there is a nonzero common
zero of $l_1,l_2$ because two homogeneous linear equations in three
unknowns have a nonzero solution. At that point every term in
the product-rule gradient vanishes. Thus such a product has a
nonzero singular point, unlike $p$. Since every element in the
$GL_3$ orbit of $xyz$ is such a product, this proves ordinary
orbit nonmembership for the source pair.

This argument alone does not compute any coordinate-ring multiplicity,
and it has not supplied the algebraic closedness or a separating
polynomial needed to promote orbit nonmembership to the exact orbit-
closure multiplicity claim. The latter separation is known in the
primary source, but is not newly proved by this gradient calculation.
For another boundary check, the analogous binary degree-two power
sum is already a product:
$x^2+y^2=(x+iy)(x-iy)$. The two linear forms are independent,
so this power sum lies in the $GL_2$ orbit of $xy$.
The auxiliary rank-one quadratic pair above therefore must not be
confused with the source's degree-two power-sum/product pair.

\section{Why symmetry data for an orbit need an extension argument}
Consider the elementary action of $\mathbb C^*$ on the affine line
by multiplication. The orbit of 1 is $\mathbb C^*$ and its closure
is $\mathbb C$. Their coordinate rings are respectively
$\mathbb C[t,t^{-1}]$ and $\mathbb C[t]$: on the punctured line,
inverting $t$ is allowed, whereas a polynomial regular on the full
line cannot have a negative power. For every integer $r$, the monomial
$t^r$ has coordinate-action character $a\mapsto a^{-r}$ on the orbit;
only $r\ge0$ occur in the closure's coordinate ring. In particular
the orbit's vector $t^{-1}$ does not extend across zero.

The injection of a closure's coordinate ring into that of its dense
orbit therefore provides upper bounds from orbit multiplicities, not
automatic lower bounds in the closure. A strategy comparing symmetry-
group invariant dimensions for two orbits needs an extension argument
on the side where a closure lower bound is required. This example
identifies the precise missing implication; it does not refute the
special extension theorems proved in the cited GCT work.

\section{Model-assisted audit and remaining scope}
The audit checked coordinate action versus partition labels, exact
multiplicities in the quotient, the difference between the auxiliary
rank-one pair and the actual power-sum/product pair, and the direction
of the orbit-ring inclusion. The only computed multiplicities here
are proved algebraically above; no symbolic package, numerical
experiment, or theorem prover was used. The original family is
underspecified and the hard permanent/determinant separation is not
obtained. Known source-specific obstructions are not claimed as new,
and no completion credit is assigned for this illustrative rederivation.

\section{Process and contributions}
Mathematics and model-assisted adversarial self-audit: OpenAI gpt-6-astra, ultra reasoning, as dispatched for this research batch; exact serving revision not exposed. Codex (GPT-6; exact serving revision not exposed) materially assisted with orchestration and format review. No human mathematical authorship, mathematical edits, or independent human proof checking is asserted.
The coordinating Codex (GPT-6; exact serving revision not exposed) also provided a separate model-assisted written-proof review of this manuscript. This is not independent human verification.
Contribution prompt A was applied to the full definition and supplied source row, then prompt B was applied adversarially to the derived results. The item-specific audit above records the scope checks and rejected extrapolations. Local assembly and structural checks do not establish mathematical correctness. The source catalogue status is preserved in TOP\_PROBLEM and is separate from this attempt status.
\paragraph{Integration review.} Codex (GPT-6; exact serving revision not exposed) reviewed the written arguments and expanded layout and source-reference wrappers. Repository metadata and the production web renderer were checked using Python 3.13 and JavaScript with MathJax 3.2.2. This model-assisted review is not independent human proof verification; no PDF compilation or formal proof checking is asserted.

\section{Completion Estimate}
\noindent\textbf{COMPLETION: 0\%}
The proved material is a rederivation of known or auxiliary facts rather than a substantiated new advance on the original target. The explicit multiplicity calculation concerns a labeled auxiliary binary-quadratic pair. The source power-sum/product calculation separates ordinary orbits only, and the missing orbit-to-closure extension prevents a new obstruction for the intended GCT targets.
\section{References}
The complete original Sources subsection above is retained, including catalogue attribution and licence notices. The additional list records primary sources actually consulted for this attempt on 11 October 2026; uncited source links reproduced above were not all rechecked.
\begin{itemize}
\item Ikenmeyer, joint work with Kandasamy, OWR 9/2020, pp. 557--560; source orbit/closure distinction, permanent/determinant motivation and power-sum/product family read in the cached primary report. \url{https://doi.org/10.4171/owr/2020/9}.
\item Ikenmeyer and Kandasamy, Implementing geometric complexity theory: On the separation of orbit closures via symmetries (2019); primary abstract inspected, full multiplicity-obstruction proof not audited. \url{https://arxiv.org/abs/1911.03990}.
\end{itemize}
\section{Final status}
\textbf{Unresolved in this attempt.} The explicit multiplicity calculation concerns a labeled auxiliary binary-quadratic pair. The source power-sum/product calculation separates ordinary orbits only, and the missing orbit-to-closure extension prevents a new obstruction for the intended GCT targets.
\end{document}
